\documentclass[11pt,a4paper]{article}
\usepackage[utf8]{inputenc}
\usepackage[T1]{fontenc}
\usepackage[french]{babel}
\usepackage{geometry}
\usepackage{xcolor}
\usepackage{hyperref}
\usepackage{titlesec}

\geometry{margin=1in}
\definecolor{titleblue}{RGB}{0,51,102}

\hypersetup{
    colorlinks=true,
    linkcolor=titleblue,
    urlcolor=titleblue,
    pdftitle={Lettre de Motivation - Ismail AISSAOUI},
}

\begin{document}

\noindent \textbf{Ismail AISSAOUI} \\
Ingénieur d'État en Intelligence Artificielle (ESI-SBA) \\
Email: ismail.aissaoui.pro@gmail.com \\
Téléphone: +213 660 70 77 96 \\
LinkedIn: https://ismail-aissaoui.vercel.app

\bigskip

\begin{flushright}
    Au Comité de Sélection du Doctorat \\
    \textbf{Programme EDT - Université de Pau (UPPA)} \\
    À l'attention de : Dr. Johann Bourcier \\
    \medskip
    1er mai 2026
\end{flushright}

\bigskip

\noindent \textbf{Objet : Candidature au doctorat : Exploration intelligente de scénarios pour les Jumeaux Numériques via les enveloppes de validité (Rang 4 - PC1)}

\bigskip

Monsieur le Docteur,

Ingénieur d'État en Intelligence Artificielle diplômé de l'ESI-SBA, je vous adresse avec enthousiasme ma candidature pour le poste de doctorat intitulé « Exploration intelligente de scénarios pour les Jumeaux Numériques via les enveloppes de validité » au sein du programme national Engineering Digital Twin (EDT). Cette candidature est motivée par la forte adéquation de mon profil avec les objectifs scientifiques de votre projet.

Actuellement en fin de cycle ingénieur chez Sonatrach, je développe un **Jumeau Numérique informé par la physique** pour la surveillance de turbines à gaz. Une partie centrale de mon travail consiste à créer des modèles génératifs permettant aux ingénieurs d'explorer des scénarios « what-if » tout en garantissant qu'ils restent dans les enveloppes de validité physique et statistique du modèle. J'ai acquis une solide expérience dans l'utilisation des architectures **Mamba-SSM** et **xLSTM** pour simuler des trajectoires temporelles complexes sous contraintes d'incertitude.

Votre projet m'intéresse vivement car il s'attaque à la question fondamentale de la confiance dans les simulations à grande échelle (territoires, environnement). Je suis particulièrement motivé par l'idée de formaliser les « enveloppes de validité » pour guider l'exploration de scénarios, en utilisant ma maîtrise de la quantification de l'incertitude (UQ) et de l'apprentissage automatique scientifique (Scientific ML).

Intégrer l'UPPA et le consortium EDT me permettrait de mettre mes compétences au service d'enjeux territoriaux majeurs tout en contribuant au développement de la **plateforme Artemis**.

Je vous remercie de l'attention que vous porterez à ma candidature et reste à votre entière disposition pour un entretien.

Dans l'attente de votre réponse, je vous prie d'agréer, Monsieur le Docteur, l'expression de mes salutations distinguées.

\bigskip

\noindent Ismail AISSAOUI

\end{document}
